\documentclass[11pt,a4paper]{article}
\usepackage{fontspec}
\IfFontExistsTF{Times New Roman}{\setmainfont{Times New Roman}}{\setmainfont{Tinos}}
\IfFontExistsTF{Consolas}{\setmonofont{Consolas}[Scale=0.90]}{\setmonofont{Liberation Mono}[Scale=0.90]}
\usepackage[english]{babel}
\usepackage[top=1.35cm,bottom=1.35cm,left=1.65cm,right=1.65cm]{geometry}
\usepackage{xcolor,enumitem,listings,booktabs,array,amsmath,amssymb,titlesec,fancyhdr,parskip,tabularx,lastpage,hyperref}
\usepackage{tcolorbox}
\tcbuselibrary{skins}
\definecolor{primaryblue}{RGB}{18,65,115}
\definecolor{secondaryteal}{RGB}{14,102,110}
\definecolor{CodeBg}{RGB}{248,250,252}
\definecolor{DarkGray}{RGB}{55,55,55}
\definecolor{boxbg}{RGB}{248,250,252}
\setlength{\parindent}{0pt}\setlength{\parskip}{0pt}
\setlist[itemize]{leftmargin=1.3em,itemsep=1pt,topsep=1pt,parsep=0pt,partopsep=0pt}
\setlist[enumerate]{leftmargin=1.4em,itemsep=1pt,topsep=1pt,parsep=0pt,partopsep=0pt}
\setlist[enumerate,2]{topsep=0.5pt,itemsep=0pt,parsep=0pt,partopsep=0pt}
\lstset{language=Python,basicstyle=\ttfamily\scriptsize,keywordstyle=\color{blue!75!black}\bfseries,stringstyle=\color{red!65!black},commentstyle=\color{primaryblue!80!black}\itshape,showstringspaces=false,columns=fullflexible,keepspaces=true,breaklines=true,frame=single,rulecolor=\color{primaryblue!45},backgroundcolor=\color{CodeBg},xleftmargin=0.2em,xrightmargin=0.2em,aboveskip=0.2em,belowskip=0.2em,tabsize=4}
\pagestyle{fancy}\fancyhf{}\setlength{\headheight}{14pt}
\fancyhead[L]{\footnotesize\textbf{DSAI1003 -- Fundamental Programming Concepts in Python}}
\fancyhead[R]{\footnotesize\textbf{Objects \& Classes Exam -- 60 Minutes}}
\fancyfoot[L]{\scriptsize Faculty of Data Science and Artificial Intelligence -- National Economics University}
\fancyfoot[R]{\footnotesize Page \thepage\ of \pageref{LastPage}}
\renewcommand{\headrulewidth}{0.4pt}\renewcommand{\footrulewidth}{0.3pt}
\titleformat{\section}{\large\bfseries\color{primaryblue}}{}{0pt}{}
\titleformat{\subsection}{\normalsize\bfseries\color{primaryblue}}{}{0pt}{}
\titlespacing*{\section}{0pt}{3pt}{1pt}\titlespacing*{\subsection}{0pt}{2pt}{1pt}
\newtcolorbox{headerbox}{colback=white,colframe=primaryblue,boxrule=0.7pt,arc=1.5mm,left=2.5mm,right=2.5mm,top=1mm,bottom=1mm}
\newtcolorbox{answerbox}[2][]{enhanced,colback=white,colframe=DarkGray!45,boxrule=0.5pt,arc=1mm,left=2mm,right=2mm,top=1mm,bottom=1mm,title=\footnotesize\color{DarkGray}\textbf{#2},#1}
\hypersetup{colorlinks=true,linkcolor=primaryblue,urlcolor=primaryblue,pdftitle={DSAI1003 - Objects and Classes Progress Exam},pdfauthor={Faculty of Data Science and Artificial Intelligence - NEU}}
\begin{document}
% PAGE 1: institutional header, answer sheet, Q1-Q5
\noindent
\begin{minipage}[t]{0.63\textwidth}\raggedright\small
\textbf{NATIONAL ECONOMICS UNIVERSITY}\\
\textbf{COLLEGE OF TECHNOLOGY}\\
\textbf{FACULTY OF DATA SCIENCE \& ARTIFICIAL INTELLIGENCE (FDA)}
\end{minipage}\hfill
\begin{minipage}[t]{0.35\textwidth}\raggedleft\footnotesize
\textbf{PROGRESS EXAMINATION}\\
Academic Year: 2026--2027\\
Semester: Fall 2026 \quad \textbf{DSAI1003}
\end{minipage}
\vspace{0.1em}
\begin{center}
{\Large\bfseries\color{primaryblue} FUNDAMENTAL PROGRAMMING CONCEPTS IN PYTHON}\\[0.05em]
{\normalsize\bfseries PROGRESS EXAM: OBJECTS AND CLASSES (CHAPTER 9)}\\[0.05em]
{\normalsize\textbf{Paper Code:} \rule{1.5cm}{0.4pt} \quad | \quad \textbf{Duration: 60 Minutes} (Excluding distribution)}
\end{center}
\begin{headerbox}
\begin{tabularx}{\linewidth}{@{}X X@{}}
\textbf{Candidate Full Name:} \dotfill & \textbf{Student ID (MSSV):} \dotfill \\[0.15em]
\textbf{Class / Cohort:} \dotfill & \textbf{Exam Room / Seat No.:} \dotfill
\end{tabularx}
\vspace{1.5mm}
\centering\setlength{\tabcolsep}{4pt}\renewcommand{\arraystretch}{1.05}
\begin{tabular}{|l|c|c|c|c||c|c|}\hline
\textbf{Section} & \textbf{A} & \textbf{B} & \textbf{C} & \textbf{D} & \textbf{Total (100)} & \textbf{Examiner Signature} \\ \hline
\textbf{Max Score} & 20 & 25 & 25 & 30 & \textbf{100} & \\[0.2em] \cline{1-6}
\textbf{Score} & & & & & & \\[0.35em] \hline
\end{tabular}
\vspace{1mm}\raggedright\scriptsize
\textbf{Instructions:} Closed-book examination; no electronic devices or Internet access. Answer all questions. Unless otherwise specified, assume Python 3 and standard runtime semantics. Standard built-in libraries (such as \texttt{math}) are permitted; third-party packages (such as NumPy, Pandas) are strictly prohibited. For program-design tasks, implement modular class interfaces with clean encapsulation and a \texttt{main()} demonstration driver. Do not detach pages.
\end{headerbox}
\vspace{1mm}
\textbf{MULTIPLE-CHOICE ANSWER SHEET} \textit{(Write A, B, C, or D in each cell.)}
\begin{center}\setlength{\tabcolsep}{5pt}\renewcommand{\arraystretch}{1.1}
\begin{tabular}{|c|c|c|c|c|c|c|c|c|c|c|}\hline
\textbf{Question} & \textbf{1} & \textbf{2} & \textbf{3} & \textbf{4} & \textbf{5} & \textbf{6} & \textbf{7} & \textbf{8} & \textbf{9} & \textbf{10} \\ \hline
\textbf{Answer} & & & & & & & & & & \\[1.5mm] \hline
\end{tabular}\end{center}
\vspace{-2.5mm}
\section*{Section A: Object-Oriented Programming Foundations \hfill \normalsize\normalfont(20 points)}
\textit{Answer all 10 questions. Each correct answer is worth 2 points. Suggested time: 12 minutes.}
\begin{enumerate}[label=\textbf{Q\arabic*.},leftmargin=2em,itemsep=1.5pt,topsep=1pt]
\item What is the relationship between a \textbf{class} and an \textbf{object} in Python?
\begin{enumerate}[label=\Alph*.,leftmargin=2em]
\item A class is an execution-time instance of an object stored in cache memory.
\item A class is a blueprint or template specifying state and behavior; an object is a concrete instance of that class.
\item A class contains only executable functions, whereas an object contains only raw data fields without methods.
\item Classes and objects are completely synonymous terms with identical memory footprints in Python.
\end{enumerate}
\item What is the primary role of the special method \texttt{\_\_init\_\_} in Python class design?
\begin{enumerate}[label=\Alph*.,leftmargin=2em]
\item It allocates hardware RAM from the operating system before the object can be created.
\item It acts as the constructor method that automatically initializes the instance variables of a newly created object.
\item It destroys the object and frees its heap space once the object goes out of scope.
\item It must be called explicitly by client code (e.g., \texttt{obj.\_\_init\_\_()}) before calling any accessor method.
\end{enumerate}
\item Why must instance methods in Python include \texttt{self} as their first parameter in the method signature?
\begin{enumerate}[label=\Alph*.,leftmargin=2em]
\item It is a reserved keyword that binds the method to the Python operating system kernel.
\item It provides an explicit reference to the specific receiver instance whose attributes are being accessed or modified.
\item It forces the method to run asynchronously in a separate worker thread.
\item The caller is required to supply \texttt{self} as an explicit argument during every method invocation.
\end{enumerate}
\item In Python conventions, what does a leading underscore in an attribute name (e.g., \texttt{self.\_balance}) signify?
\begin{enumerate}[label=\Alph*.,leftmargin=2em]
\item The Python virtual machine strictly prevents external code from reading or modifying the attribute.
\item It signals an internal implementation detail (encapsulation) that external users should not access directly.
\item It marks the variable as a global constant that cannot be assigned new values.
\item It automatically generates compiled getter and setter functions at runtime.
\end{enumerate}
\item In the \texttt{CashRegister} class from the textbook (\texttt{addItem(price)}, \texttt{getTotal()}, \texttt{getCount()}, \texttt{clear()}), which pair correctly classifies an \textbf{accessor} method and a \textbf{mutator} method?
\begin{enumerate}[label=\Alph*.,leftmargin=2em]
\item Accessor: \texttt{addItem()}; Mutator: \texttt{getTotal()}
\item Accessor: \texttt{clear()}; Mutator: \texttt{getCount()}
\item Accessor: \texttt{getTotal()}; Mutator: \texttt{addItem()}
\item Accessor: \texttt{getCount()}; Mutator: \texttt{getTotal()}
\end{enumerate}
\end{enumerate}
\newpage
% PAGE 2: Q6-Q10 and Q11
\begin{enumerate}[label=\textbf{Q\arabic*.},leftmargin=2em,itemsep=1.5pt,topsep=1pt,resume]
\item Consider the statements:\newline\texttt{c1 = Counter(); c1.click(); c2 = c1; c2.click()}\newline What are the values returned by \texttt{c1.getValue()} and \texttt{c2.getValue()}?
\begin{enumerate}[label=\Alph*.,leftmargin=2em]
\item \texttt{c1.getValue()} is 1; \texttt{c2.getValue()} is 2
\item \texttt{c1.getValue()} is 2; \texttt{c2.getValue()} is 2, because \texttt{c1} and \texttt{c2} refer to the exact same object
\item \texttt{c1.getValue()} is 1; \texttt{c2.getValue()} is 1
\item A \texttt{ReferenceError} is raised due to duplicate object binding
\end{enumerate}
\item Given two distinct instances: \texttt{acc1 = BankAccount(500.0)} and \texttt{acc2 = BankAccount(500.0)}. What does \texttt{acc1 is acc2} evaluate to?
\begin{enumerate}[label=\Alph*.,leftmargin=2em]
\item \texttt{True}, because both instances store identical internal balance numbers.
\item \texttt{True}, because both instances belong to the identical \texttt{BankAccount} class.
\item \texttt{False}, because \texttt{is} tests object identity (memory address), and \texttt{acc1} and \texttt{acc2} are separate objects.
\item A \texttt{TypeError} is raised because user-defined classes do not support identity comparisons.
\end{enumerate}
\item Which statement correctly contrasts \textbf{instance variables} and \textbf{class variables}?
\begin{enumerate}[label=\Alph*.,leftmargin=2em]
\item Instance variables are bound to individual objects via \texttt{self}; class variables are shared across all instances.
\item Class variables are mutable, whereas instance variables are strictly immutable in Python.
\item Instance variables are declared outside all methods, whereas class variables are declared inside \texttt{\_\_init\_\_}.
\item Class variables are stored on the call stack, while instance variables are stored in CPU registers.
\end{enumerate}
\item When calling \texttt{print(obj)} or \texttt{str(obj)}, which special method is automatically invoked if implemented by the class?
\begin{enumerate}[label=\Alph*.,leftmargin=2em]
\item \texttt{\_\_init\_\_} \quad \Alph*.~ \texttt{\_\_display\_\_} \quad \Alph*.~ \texttt{\_\_str\_\_} \quad \Alph*.~ \texttt{\_\_format\_out\_\_}
\end{enumerate}
\item If variable \texttt{active\_card = None}, what occurs when executing \texttt{active\_card.tap\_ride("Central", 2.50)}?
\begin{enumerate}[label=\Alph*.,leftmargin=2em]
\item The statement executes silently with no operation performed.
\item A new default instance of the expected class is automatically constructed.
\item An \texttt{AttributeError} is raised because \texttt{'NoneType'} has no attribute or method named \texttt{'tap\_ride'}.
\item Python converts \texttt{None} to a boolean \texttt{False} and returns \texttt{0.0}.
\end{enumerate}
\end{enumerate}
\vspace{-1mm}
\section*{Section B: Output Tracing, Object Mechanics \& Diagnostics \hfill \normalsize\normalfont(25 points)}
\textit{Answer Q11--Q15. Each question is worth 5 points. Suggested time: 15 minutes.}
\subsection*{Q11. Exact Output of Multi-Object State Tracing (5 points)}
Predict the \textbf{exact console output} of the following Python program, including spacing and punctuation.
\begin{lstlisting}
class SmartCounter:
    def __init__(self, step=1):
        self._count = 0
        self._step = step
    def click(self): self._count += self._step
    def boost(self, factor): self._step *= factor
    def get_value(self): return self._count

c1 = SmartCounter(2)
c2 = SmartCounter(5)
c3 = c1
c1.click()
c2.click()
c3.boost(3)
c3.click()
c2.click()
print(f"c1: {c1.get_value()} | c2: {c2.get_value()}")
print(f"c1 is c3: {c1 is c3} | c1 is c2: {c1 is c2}")
\end{lstlisting}
\begin{answerbox}[height=2.7cm]{Output Box for Q11}\end{answerbox}
\newpage
% PAGE 3: Q12-Q15
\subsection*{Q12. Evaluate Class and Object Expressions (5 points)}
Write the exact evaluated Python result or raised exception class for each independent expression (1 point each):
\begin{enumerate}[label=\alph*),leftmargin=2em,itemsep=2.5pt]
\item Given \texttt{class Item: pass; it = Item()}, the value of \texttt{isinstance(it, Item)}: \hrulefill
\item The exception raised by executing: \lstinline|class T: pass; obj = T(); print(obj.price)|: \hrulefill
\item After \lstinline|a = [10, 20]; b = a; b.append(30)|, the value of \texttt{len(a)}: \hrulefill
\item In a reduced \texttt{Fraction(4, 12)} class, the value stored in attribute \texttt{\_numerator}: \hrulefill
\item The evaluated boolean value of \texttt{bool(None)}: \hrulefill
\end{enumerate}
\subsection*{Q13. Diagnose and Correct Three Object-Oriented Defects (5 points)}
For each defect, identify the underlying bug and provide the corrected Python code.\par
\renewcommand{\arraystretch}{1.20}\small
\begin{tabularx}{\linewidth}{|p{4.8cm}|X|X|}\hline
\textbf{Erroneous Code} & \textbf{Problem / Cause} & \textbf{Correction} \\ \hline
{\ttfamily\scriptsize class Account:}\newline{\ttfamily\scriptsize \ \ def deposit(amount):}\newline{\ttfamily\scriptsize \ \ \ \ self.\_balance += amount}\newline{\itshape Goal: add funds to balance} & & \\[6.5mm]\hline
{\ttfamily\scriptsize class Counter:}\newline{\ttfamily\scriptsize \ \ def \_\_init\_\_(self): self.\_val = 0}\newline{\ttfamily\scriptsize \ \ def click(self): \_val += 1}\newline{\itshape Goal: increment counter} & & \\[6.5mm]\hline
{\ttfamily\scriptsize class Student:}\newline{\ttfamily\scriptsize \ \ def \_\_init\_\_(self, name, tags=[]):}\newline{\ttfamily\scriptsize \ \ \ \ self.\_tags = tags}\newline{\itshape Goal: assign independent list} & & \\[6.5mm]\hline
\end{tabularx}\normalsize
\subsection*{Q14. Bank Account State Tracing with Overdraft \& Fee Rules (5 points)}
Predict the exact console output of the code below and show the balance after each method invocation.
\begin{lstlisting}
class BankAccount:
    def __init__(self, initial=0.0): self._bal = initial
    def deposit(self, amt): self._bal += amt
    def withdraw(self, amt):
        if amt > self._bal: self._bal -= 10.0   # $10 overdraft penalty fee
        else: self._bal -= amt
    def add_interest(self, rate): self._bal += self._bal * (rate / 100.0)
    def get_balance(self): return self._bal

acc = BankAccount(200.0)
acc.deposit(100.0); acc.withdraw(350.0); acc.withdraw(50.0); acc.add_interest(5.0)
print(f"Final Balance: {acc.get_balance():.2f}")
\end{lstlisting}
\begin{answerbox}[height=2.3cm]{Output and Trace Table for Q14}\end{answerbox}
\subsection*{Q15. Object References vs Shallow Copies in Object Attributes (5 points)}
Analyze the code below: (1) Output of \texttt{c1.items()}; (2) Output of \texttt{c2.items()}; (3) Why modifying \texttt{raw} impacted \texttt{c1} but did NOT impact \texttt{c2}.
\begin{lstlisting}
class Cart:
    def __init__(self, items): self._items = items
    def add(self, it): self._items.append(it)
    def items(self): return self._items

raw = ["Book", "Pen"]
c1 = Cart(raw); c2 = Cart(list(raw))
raw.append("Ruler"); c1.add("Eraser")
\end{lstlisting}
\begin{answerbox}[height=2.1cm]{Answers and Technical Explanation for Q15}\end{answerbox}
\newpage
% PAGE 4: design
\section*{Section C: Object-Oriented Design \& Modular Class Decomposition \hfill \normalsize\normalfont(25 points)}
\textit{Answer all four tasks. Suggested time: 15 minutes.}
\begin{tcolorbox}[colback=boxbg,colframe=secondaryteal,arc=1.5mm,boxrule=0.6pt,top=1.5mm,bottom=1.5mm]
\textbf{Problem Specification: Smart Urban Transit Card Management System (\texttt{MetroCard})}\\[0.5mm]
\small
A municipal metro transit authority deploys an electronic stored-value fare card: \texttt{MetroCard}.
\begin{itemize}
\item \textbf{Encapsulated State:} Each card maintains a unique \texttt{card\_id} (str), a prepaid balance \texttt{\_balance} (float, $\ge 0$), cumulative fares spent today \texttt{\_daily\_spent} (float), a maximum daily fare cap \texttt{\_daily\_cap} (float, e.g., \$8.00), and a recorded list of ride logs \texttt{\_rides} containing tuples of \texttt{(station\_name, fare\_paid)}.
\item \textbf{Automated Daily Fare Capping:} Commuters pay the standard fare until cumulative spending reaches the daily cap. Once the cap is reached, subsequent rides on that day are \textbf{free} (\$0.00). If a ride's base fare would exceed the remaining cap margin, the passenger is charged \textbf{only the difference} up to the cap:
$$\text{effective\_fare} = \max\left(0.0, \, \min(\text{base\_fare}, \, \text{\_daily\_cap} - \text{\_daily\_spent})\right)$$
\item \textbf{Gate Validation:} If the passenger's balance is less than the effective fare ($\text{\_balance} < \text{effective\_fare}$), the turnstile gate denies entry: no money is deducted, no ride is logged, and the method returns \texttt{False}. Otherwise, it deducts \texttt{effective\_fare}, updates \texttt{\_daily\_spent}, appends the ride to \texttt{\_rides}, and returns \texttt{True}.
\end{itemize}
\end{tcolorbox}
\vspace{0.5mm}
\textbf{Task 1: Specify Public Methods, Parameter Types, Return Types, and Invariants (6 points)}
For each method of \texttt{MetroCard}, specify its arguments, return type, and behavioral contract:
\begin{itemize}[leftmargin=1.8em,itemsep=2pt]
\item \texttt{\_\_init\_\_(self, card\_id, initial\_balance=0.0, daily\_cap=8.0)} $\rightarrow$ Purpose: \dotfill
\item \texttt{top\_up(self, amount)} $\rightarrow$ Return Type: \rule{2.8cm}{0.4pt} \quad Precondition/Contract: \dotfill
\item \texttt{tap\_ride(self, station, base\_fare)} $\rightarrow$ Return Type: \rule{2.5cm}{0.4pt} \quad Precondition/Contract: \dotfill
\item \texttt{reset\_daily(self)} $\rightarrow$ Return Type: \rule{2.8cm}{0.4pt} \quad Postcondition: \dotfill
\end{itemize}
\textbf{Task 2: Implement the Core Fare Deduction Logic in \texttt{tap\_ride} (7 points)}
Write a Python method implementation for \texttt{tap\_ride(self, station, base\_fare)}. Correctly compute the capped fare, perform the balance check, update internal state, and return boolean status.
\begin{answerbox}[height=3.4cm]{Python Implementation for Task 2}\end{answerbox}
\textbf{Task 3: Hand-Trace a Transit Simulation Sequence (6 points)}
A student purchases a \texttt{MetroCard("MC-99", initial\_balance=15.00, daily\_cap=8.00)}. Compute the state after each tap:
\begin{enumerate}[label=\arabic*.,leftmargin=1.8em,itemsep=1pt]
\item \texttt{tap\_ride("Central", 3.50)} $\rightarrow$ Fare: \rule{1.6cm}{0.4pt} \quad Balance: \rule{1.6cm}{0.4pt} \quad Daily: \rule{1.6cm}{0.4pt}
\item \texttt{tap\_ride("WestEnd", 3.00)} $\rightarrow$ Fare: \rule{1.6cm}{0.4pt} \quad Balance: \rule{1.6cm}{0.4pt} \quad Daily: \rule{1.6cm}{0.4pt}
\item \texttt{tap\_ride("NorthGate", 3.00)} $\rightarrow$ Fare: \rule{1.6cm}{0.4pt} \quad Balance: \rule{1.6cm}{0.4pt} \quad Daily: \rule{1.6cm}{0.4pt}
\item \texttt{tap\_ride("Campus", 2.50)} $\rightarrow$ Fare: \rule{1.6cm}{0.4pt} \quad Balance: \rule{1.6cm}{0.4pt} \quad Daily: \rule{1.6cm}{0.4pt}
\end{enumerate}
Final \texttt{\_balance}: \rule{2.2cm}{0.4pt} \quad Final \texttt{\_daily\_spent}: \rule{2.2cm}{0.4pt} \quad Total successful rides: \rule{2.2cm}{0.4pt}
\vspace{0.5mm}
\textbf{Task 4: Technical Analysis \& Object-Oriented Principles (6 points)}
\begin{enumerate}[label=\alph*),leftmargin=1.8em,itemsep=1.5pt]
\item Why is encapsulation crucial in \texttt{MetroCard}? Explain what critical bugs could occur if external code could directly write to \texttt{card.\_balance = -50.0} or modify \texttt{card.\_daily\_spent}.\\[-1pt]\rule{\linewidth}{0.4pt}
\item When passing a \texttt{MetroCard} instance into a batch function \texttt{batch\_process(card)}, does Python pass a copy of the card or an object reference? What does this imply if the function mutates the card?\\[-1pt]\rule{\linewidth}{0.4pt}
\end{enumerate}
\newpage
% PAGE 5: coding problem 1
\section*{Section D: Applied Programming in Python \hfill \normalsize\normalfont(30 points)}
\textit{Write complete, modular Python 3 programs. Use clean encapsulation, defensive parameter checks, explicit \texttt{self} bindings, and a \texttt{main()} demonstration driver. Standard library only. Suggested time: 18 minutes.}
\subsection*{Problem 1: Commercial Bank Account with Tiered Overdraft Protection (15 points)}
Implement a complete \texttt{CommercialAccount} class modeling a business checking account with overdraft facilities.
\textbf{Requirements:}
\begin{enumerate}[label=\arabic*.,leftmargin=2em,itemsep=1.2pt]
\item \texttt{\_\_init\_\_(self, acc\_num, owner, balance=0.0, overdraft\_limit=100.0)}:
  \begin{itemize}
  \item Initialize attributes: \texttt{\_acc\_num} (str), \texttt{\_owner} (str), \texttt{\_balance} (float), \texttt{\_overdraft\_limit} (float), \texttt{\_ledger} (list). Raise \texttt{ValueError} if \texttt{balance < 0} or \texttt{overdraft\_limit < 0}. Log \texttt{"Account Opened"}.
  \end{itemize}
\item \texttt{deposit(self, amount)}: If \texttt{amount <= 0}, raise \texttt{ValueError}. Add \texttt{amount} to \texttt{\_balance}, append \texttt{f"Deposit: +\$\{amount:.2f\}"} to \texttt{\_ledger}, and return current balance.
\item \texttt{withdraw(self, amount)}: If \texttt{amount <= 0}, raise \texttt{ValueError}.
  \begin{itemize}
  \item If \texttt{amount <= self.\_balance}: Deduct \texttt{amount}, log \texttt{f"Standard: -\$\{amount:.2f\}"}, return \texttt{True}.
  \item Else if \texttt{amount <= self.\_balance + self.\_overdraft\_limit}: Deduct \texttt{amount + 10.0} (\$10 overdraft penalty fee), log \texttt{f"Overdraft: -\$\{amount:.2f\} (Fee: \$10.00)"}, return \texttt{True}.
  \item Else (insufficient funds): Do not change balance, log \texttt{f"Declined: -\$\{amount:.2f\}"}, return \texttt{False}.
  \end{itemize}
\item \texttt{transfer\_to(self, target\_acc, amount)}: Attempt withdrawal of \texttt{amount} from \texttt{self}. If successful, deposit into \texttt{target\_acc} and return \texttt{True}; otherwise return \texttt{False} (atomic transfer).
\item \texttt{\_\_str\_\_(self)}: Return string: \texttt{"Account[\#\{acc\_num\}, \{owner\}]: Balance = \$\{balance:.2f\}"}.
\item \texttt{main()}: Create account \#A101 with \$150.00, deposit \$50.00, execute a \$250.00 overdraft withdrawal, transfer \$30.00 to account \#B202, and print both account summaries and \#A101's audit ledger.
\end{enumerate}
\begin{answerbox}[height=8.5cm]{Python Solution for Problem 1}\end{answerbox}
\newpage
% PAGE 6: coding problem 2
\subsection*{Problem 2: Immutable Rational Number Class with Operator Overloading (\texttt{Fraction}) (15 points)}
Implement a complete, immutable \texttt{Fraction} class supporting exact rational arithmetic, directly modeling Section 9.11.
\textbf{Requirements:}
\begin{enumerate}[label=\arabic*.,leftmargin=2em,itemsep=1.2pt]
\item Constructor \texttt{\_\_init\_\_(self, numerator=0, denominator=1)}:
  \begin{itemize}
  \item If \texttt{denominator == 0}, raise \texttt{ZeroDivisionError("Denominator cannot be zero")}.
  \item If \texttt{denominator < 0}, normalize signs so the denominator is strictly positive.
  \item Reduce fraction to lowest terms by dividing by $\gcd(|\text{numerator}|, \text{denominator})$. (Import \texttt{math.gcd}).
  \item Store simplified values in \texttt{self.\_num} and \texttt{self.\_den}.
  \end{itemize}
\item Immutability: Do not provide mutator methods; all arithmetic operations return a \textbf{new} \texttt{Fraction} instance.
\item Special Dunder Methods:
  \begin{itemize}
  \item \texttt{\_\_str\_\_(self)}: Return \texttt{"0"} if \texttt{\_num == 0}; return \texttt{f"\{\_num\}"} if \texttt{\_den == 1}; else \texttt{f"\{\_num\}/\{\_den\}"}.
  \item \texttt{\_\_add\_\_(self, other)}: Return a new \texttt{Fraction} representing $\frac{n_1 d_2 + n_2 d_1}{d_1 d_2}$.
  \item \texttt{\_\_sub\_\_(self, other)}: Return a new \texttt{Fraction} representing $\frac{n_1 d_2 - n_2 d_1}{d_1 d_2}$.
  \item \texttt{\_\_mul\_\_(self, other)}: Return a new \texttt{Fraction} representing $\frac{n_1 n_2}{d_1 d_2}$.
  \item \texttt{\_\_eq\_\_(self, other)}: Return \texttt{True} if \texttt{self.\_num == other.\_num} and \texttt{self.\_den == other.\_den}.
  \end{itemize}
\item Utility Method \texttt{to\_float(self)}: Return floating-point approximation \texttt{self.\_num / self.\_den}.
\item \texttt{main()}: Instantiate $f_1 = \text{Fraction}(1, 6)$, $f_2 = \text{Fraction}(1, 3)$, demonstrate $f_1 + f_2 = 1/2$, demonstrate subtraction and multiplication, demonstrate that $\text{Fraction}(-2, -4)$ normalizes to $1/2$, and catch \texttt{ZeroDivisionError} on $\text{Fraction}(3, 0)$.
\end{enumerate}
\begin{answerbox}[height=8.8cm]{Python Solution for Problem 2}\end{answerbox}
\vfill
\begin{center}\rule{0.5\textwidth}{0.4pt}\\[1mm]
\textbf{\color{primaryblue}--- END OF EXAM PAPER ---}
\end{center}
\end{document}